\section{Real algebra: deciding by cells, not by sampling}
\label{sec:realalgebraic}

Section~\ref{sec:nonlinear} certifies that a polynomial is nonnegative. It does
not decide a quantified statement about real solutions, and the two problems
are genuinely different: a certificate search that fails proves nothing,
whereas a decision procedure must answer. Two contributing efforts supply
decision procedures over the reals, and they agree on the one thing that makes
the subject tractable and disagree about everything above it.

The agreement is on how to name a real algebraic number, which is the
representation question the whole subject turns on. One does not name
$\sqrt2$. One supplies a rational polynomial together with a rational interval
containing exactly one of its roots, and the pair \emph{is} the number. Every
downstream operation is then rational arithmetic on that data, and nothing is
approximated.

\subsection{The shared component: a complete cover of the line}

Both efforts build the same object, by the same three steps, and neither cites
the other.

Collect the polynomials whose sign can change the answer --- leading
coefficients, discriminants, resultants of pairs, and the outer problem's own
parameter polynomials. Form $P$, the monic squarefree part of their product.
Its real roots $\alpha_1<\cdots<\alpha_r$ cut the line into $2r+1$ \emph{cells}:
the $r$ points themselves, the $r-1$ open intervals between them, and the two
unbounded tails. On each cell every collected polynomial has constant sign, so
the original question has a constant answer there, and finitely many cells
decide a statement about uncountably many reals.

What makes this checkable rather than merely computable is the completeness
check, and both efforts implement the same three-part test:

\begin{enumerate}[leftmargin=1.4em]
\item each supplied interval isolates exactly one root of $P$, verified by a
  rational root count on that interval;
\item the intervals are disjoint and strictly increasing, with no endpoint a
  root of $P$;
\item their number equals the \emph{total} number of real roots of $P$.
\end{enumerate}

The third is the one that matters. Without it a certificate that quietly omits
a root describes a coarser partition on which the answer may be constant for
the wrong reason, and every cell-local check still passes. Both efforts make
the checker recompute $P$ itself from the required guards rather than accept a
supplied polynomial, for the same reason: a producer that omits a guard would
otherwise be believed.

\begin{remark}[An isolator is not an approximation]
The interval denotes the unique root inside it. Refining it changes the
representation and not the number, and a checker that compared floating-point
endpoints would be checking a different statement. One effort admits a
degenerate point descriptor $[a,a]$ for rational roots; the other forbids
finite endpoints that are roots at all. Both conventions are sound and they are
not interchangeable, so the schemas stay distinct.
\end{remark}

That shared component is exactly one module, and both efforts independently
name it as their first Lean milestone. It should be built once.

\subsection{Above the cover, they diverge completely}

\paragraph{Fibres and root indices.} One effort takes a bivariate problem and,
over each cell of the line, decomposes the fibre: it specialises the
coefficients at the cell's algebraic point, isolates the $y$-roots there, and
cuts the plane above that cell into $2n+1$ pieces. A quantified statement then
becomes a truth vector over finitely many plane cells, and --- the part that
makes it a \emph{witness} rather than a verdict --- an existential is answered
by a \emph{root index}: not a number, but the position of a root in the ordered
list over that cell. The witness is a selector, constant on each cell, and it
is checkable because the ordering is.

Determining the sign of a polynomial at an algebraic point is the technical
core, and it is done by localised fractions with a gcd zero test and interval
refinement, never by evaluation at a numerical approximation.

\paragraph{Counting without building the fibre.} The other effort never
constructs a fibre object. Over each cell it evaluates a \emph{count circuit}:
the number of real roots of the specialised system is determined from the signs
of characteristic coefficients, by Hermite's quadratic-form method and
Descartes-exact reasoning. The output per cell is an integer, and the outer
verdict follows from the integers. For univariate problems it can also return
an isolating interval, but its general answer is a count and not a witness.

These are different capabilities, not two implementations of one. A count
settles ``how many'' uniformly; a root index settles ``which one'' and can be
substituted into a source goal. Neither subsumes the other, and the merge keeps
both --- while noting that they share their entire lower half.

\subsection{Three impossibility proofs with one moral}

The most useful thing these two efforts do, and a third from the same round
does it too, is refute the natural objection that a bigger dictionary would
have sufficed. Each proves a small impossibility, by a different technique:

\begin{itemize}[leftmargin=1.4em]
\item no polynomial $p$ satisfies $p(x)^2=x^2+1$ identically;
\item no rational function $r$ satisfies $r(t)^2=t$ on an open interval;
\item no synchronised binary automaton with a fixed number of states
  recognises $y=x^2$ or $y=2^x$ --- a growth obstruction, from
  Section~\ref{sec:automatic}.
\end{itemize}

The moral is shared and the theorems are not. Each says that the relevant
\emph{representation} is the obstacle, so enlarging a search budget inside that
representation cannot help. They are three results and a provenance ledger that
cites them as one would be wrong three ways.

\subsection{The trap that only shows up after a negation}

Both efforts here, and one elsewhere in the same round, hit the same failure
mode in different clothing, and it is worth isolating because it is invisible
until someone takes a complement.

A decision procedure that reasons about the \emph{generic} case --- the open
cells, the typical parameter, the padded encoding --- can be right about every
case it considered and still return the wrong answer, because the exceptional
case it dropped was exactly where the statement failed. The real-algebraic
version is the exceptional parameter: for $xy=1$ the value $x=0$ has no
solution at all, and a procedure that reasons only over open cells will report a
universally solvable system. The automatic-structures version is the zero-tail:
an encoding that ignores trailing zeros gets $y=2x+1$ right for every
represented input and wrong at $x=0$.

In both cases the original answer is correct and the \emph{negation} is false,
which is why testing the positive direction finds nothing. The discipline both
efforts adopt is the same: the exceptional cells are cells, carried through the
whole computation and checked, not special-cased at the end.

\subsection{What these decisions cost, and what they do not claim}

The cell count is the cost, and it is not small. One effort's corpus tops out
at a cover polynomial of degree 25 and fifteen cells for a quintic with a
threshold; the other records 356 plane cells across 24 examples. Neither claims
scalability, and one states the position plainly: these are exact procedures on
small inputs, and a general cylindrical decomposition is doubly exponential no
matter how the cells are certified.

Two further scope conditions are worth carrying. A count is not a witness, and
a witness index is not a closed form --- substituting a root index into a Lean
goal requires the ordering theorem, not just the index. And both producers use
the same computer-algebra backend for their untrusted search; since both also
share this round's environment, a defect in that backend would be a
\emph{common-mode} risk across two proposals rather than one, which neither
effort is in a position to notice about the other and which this document is.
